\clearpage\SourceLeaf{124}%
\pdfbookmark[1]{The Zodiac of Juno}{zodiacus1805}
\begin{center}
{\small THE ZODIAC OF JUNO.\par}
\medskip
{\small (A supplement to the paper in the August 1804 issue of \textit{M. C.} [p.\,106 of this volume])\par}
\medskip\hrule\medskip
{\small Monatliche Correspondenz zur Beförderung der Erd- und Himmels-Kunde,\par
edited by Baron \textsc{von Zach}. March 1805.\par}
\medskip\hrule
\end{center}\medskip
Although Juno has been observed for only a short time, the elements of
its orbit already seem accurate enough to mark out the limits of the
zone in which the planet can appear to us, for the atlas one may hope
to see from the planet’s distinguished discoverer. At the request of my friend
\textsc{Harding}, I was all the more glad to undertake this work
because small, faint Juno, in positions less favourable than those
it occupied this year, will be the hardest of the three new planets
to observe and will therefore have the greatest need for detailed
star charts.

According to my latest observations, my fourth set of elements for
Juno still seems to agree so well with its motion that I would be
unable to make a reliable further improvement to the orbit at this
point. I have therefore used these elements as the basis of my
calculations and taken the following values in the notation of the
paper mentioned above.
\[
\begin{aligned}
e&=0.25684\\
k&=2.49630\\
g&=62^0\;19'\;35''\\
e'&=0.01679\\
k'&=0.999718\\
g'&=108^0\;30'\;1''
\end{aligned}
\]
\SourceLeaf{125}%
From this I found the following equation for the condition.
\[
\cos t'=0.400477\cos t-0.053103
\]

The values of $\cos t'$ therefore lie between $+0.347374$ and
$-0.453580$. The values of $t'$ consequently lie between
$69^0\;40'\;23''$ and $116^0\;58'\;25''$ on the one hand, and between
$234^0\;5'\;47''$ and $281^0\;23'\;49''$ on the other. The heliocentric
positions of the Earth from which Juno can appear at its limits thus
lie between longitudes $240^0\;45'$ and $288^0\;3'$, and between
$54^0\;6'$ and $101^0\;24'$. The former occur roughly from 22\,May to
10\,July, when Juno must appear once at its northern limit. The
latter occur from 16\,November to 2\,January, when it reaches its
southern limit once.

On this basis, just as for Ceres and Pallas, $36$ points were
determined on the northern limit and the same number on the southern
limit. The following formulas and constants are used to calculate
the right ascensions and declinations directly.
\[
\begin{array}{rcrrr}
a&=&87^0&59'&22''\\
b&=&100&32&38\\
c&=&10&44&17\\
A&=&261&18&7\\
B&=&171&40&35\\
C&=&160&38&2
\end{array}
\]
\[
v=\text{true anomaly of Juno}
\]
\[
\begin{aligned}
x&=\frac{\alpha\sin(v+323^0\;37'\;41'')}{1-e\cos v}\\[8pt]
y&=\frac{\ObservedBeta\sin(v+234^0\;0'\;9'')}{1-e\cos v}\\[8pt]
z&=\frac{\gamma\sin(v+222^0\;57'\;37'')}{1-e\cos v}
\end{aligned}
\]
\[
\begin{aligned}
\alpha&=k\sin a; &\log\alpha&=0.397030\\
\ObservedBeta&=k\sin b; &\log\ObservedBeta&=0.389902\\
\gamma&=k\sin c; &\log\gamma&=9.667556
\end{aligned}
\]

Between these $72$ points, $144$ others were inserted. Their
right ascensions increase from one multiple of $5$ degrees to the
next, in steps of $5$ degrees. They are shown in the following table.
\SourceLeaf{126}
\begin{center}
\begingroup\small\setlength{\tabcolsep}{3pt}
\begin{tabular}{r|rrl|rrl||r|rrl|rrl}
\multicolumn{14}{c}{\textit{Zodiac} of Juno.}\\[4pt]
&\multicolumn{6}{c||}{Declination}&&\multicolumn{6}{c}{Declination}\\
R.A.&\multicolumn{3}{c|}{\shortstack{at northern\\limit}}&\multicolumn{3}{c||}{\shortstack{at southern\\limit}}&R.A.&\multicolumn{3}{c|}{\shortstack{at northern\\limit}}&\multicolumn{3}{c}{\shortstack{at southern\\limit}}\\\hline
$0^0$&$4^0$&$23'$&N.&$10^0$&$38'$&S.&$180^0$&$6^0$&$59'$&N.&$3^0$&$9'$&S.\\
$5$&$5$&$13$&&$10$&$8$&&$185$&$6$&$5$&&$3$&$47$&\\
$10$&$6$&$3$&&$9$&$36$&&$190$&$5$&$11$&&$4$&$27$&\\
$15$&$6$&$51$&&$9$&$3$&&$195$&$4$&$17$&&$5$&$8$&\\
$20$&$7$&$39$&&$8$&$29$&&$200$&$3$&$23$&&$5$&$50$&\\
$25$&$8$&$26$&&$7$&$54$&&$205$&$2$&$30$&&$6$&$32$&\\
$30$&$9$&$11$&&$7$&$19$&&$210$&$1$&$38$&&$7$&$15$&\\
$35$&$9$&$54$&&$6$&$43$&&$215$&$0$&$49$&&$7$&$58$&\\
$40$&$10$&$34$&&$6$&$7$&&$220$&$0$&$11$&&$8$&$40$&\\
$45$&$11$&$13$&&$5$&$32$&&$225$&$0$&$44$&S.&$9$&$21$&\\
$50$&$11$&$49$&&$4$&$56$&&$230$&$1$&$26$&&$10$&$0$&\\
$55$&$12$&$23$&&$4$&$21$&&$235$&$2$&$3$&&$10$&$38$&\\
$60$&$12$&$53$&&$3$&$47$&&$240$&$2$&$38$&&$11$&$13$&\\
$65$&$13$&$21$&&$3$&$14$&&$245$&$3$&$8$&&$11$&$45$&\\
$70$&$13$&$45$&&$2$&$43$&&$250$&$3$&$34$&&$12$&$15$&\\
$75$&$14$&$7$&&$2$&$13$&&$255$&$3$&$55$&&$12$&$42$&\\
$80$&$14$&$25$&&$1$&$45$&&$260$&$4$&$12$&&$13$&$5$&\\
$85$&$14$&$40$&&$1$&$19$&&$265$&$4$&$24$&&$13$&$26$&\\
$90$&$14$&$51$&&$0$&$56$&&$270$&$4$&$32$&&$13$&$43$&\\
$95$&$14$&$58$&&$0$&$35$&&$275$&$4$&$34$&&$13$&$56$&\\
$100$&$15$&$1$&&$0$&$17$&&$280$&$4$&$32$&&$14$&$7$&\\
$105$&$15$&$1$&&$0$&$2$&&$285$&$4$&$25$&&$14$&$14$&\\
$110$&$14$&$55$&&$0$&$10$&N.&$290$&$4$&$14$&&$14$&$17$&\\
$115$&$14$&$46$&&$0$&$19$&&$295$&$3$&$58$&&$14$&$18$&\\
$120$&$14$&$32$&&$0$&$24$&&$300$&$3$&$38$&&$14$&$16$&\\
$125$&$14$&$14$&&$0$&$26$&&$305$&$3$&$14$&&$14$&$11$&\\
$130$&$13$&$52$&&$0$&$24$&&$310$&$2$&$45$&&$14$&$2$&\\
$135$&$13$&$26$&&$0$&$19$&&$315$&$2$&$14$&&$13$&$52$&\\
$140$&$12$&$56$&&$0$&$9$&&$320$&$1$&$38$&&$13$&$39$&\\
$145$&$12$&$22$&&$0$&$4$&S.&$325$&$1$&$0$&&$13$&$24$&\\
$150$&$11$&$44$&&$0$&$20$&&$330$&$0$&$19$&&$13$&$6$&\\
$155$&$11$&$3$&&$0$&$40$&&$335$&$0$&$25$&N.&$12$&$46$&\\
$160$&$10$&$19$&&$1$&$4$&&$340$&$1$&$9$&&$12$&$24$&\\
$165$&$9$&$32$&&$1$&$31$&&$345$&$1$&$57$&&$12$&$0$&\\
$170$&$8$&$43$&&$2$&$1$&&$350$&$2$&$45$&&$11$&$35$&\\
$175$&$7$&$52$&&$2$&$34$&&$355$&$3$&$34$&&$11$&$7$&\\
$180$&$6$&$59$&&$3$&$9$&&$360$&$4$&$23$&&$10$&$38$&\\
\end{tabular}
\endgroup
\end{center}
\begin{center}\rule{.25\linewidth}{.4pt}\end{center}
